\documentclass[11pt]{article}
\usepackage[T1]{fontenc}
\usepackage[utf8]{inputenc}
\usepackage{lmodern}
\usepackage[a4paper,margin=27mm]{geometry}
\usepackage{amsmath,amssymb,amsthm,mathtools,mathrsfs}
\usepackage{microtype}
\usepackage{booktabs,array,longtable,enumitem}
\usepackage{xcolor}
\usepackage[hidelinks]{hyperref}
\usepackage{bookmark}
\hypersetup{
 pdftitle={Erdos-Straus Selectors at the Koide-S6 Determinant Boundary},
 pdfauthor={Technical research continuation}
}
\emergencystretch=2em
\setlist{nosep,leftmargin=2.1em}

\newtheorem{theorem}{Theorem}[section]
\newtheorem{proposition}[theorem]{Proposition}
\newtheorem{lemma}[theorem]{Lemma}
\newtheorem{corollary}[theorem]{Corollary}
\theoremstyle{definition}
\newtheorem{definition}[theorem]{Definition}
\theoremstyle{remark}
\newtheorem{remark}[theorem]{Remark}

\newcommand{\Z}{\mathbb Z}
\newcommand{\Q}{\mathbb Q}
\newcommand{\R}{\mathbb R}
\newcommand{\Kfield}{\mathbb Q(\sqrt2,\sqrt3)}
\newcommand{\Sym}{\operatorname{Sym}}
\newcommand{\diag}{\operatorname{diag}}
\newcommand{\coker}{\operatorname{coker}}
\newcommand{\SNF}{\operatorname{SNF}}
\newcommand{\ES}{\operatorname{ES}}
\newcommand{\im}{\operatorname{Im}}
\newcommand{\rk}{\operatorname{rank}}
\newcommand{\dd}{\,\mathrm d}

\title{Erd\H{o}s--Straus Selectors at the Koide--$S^6$ Determinant Boundary\\
\large Exact incidence maps, primitive rays, terminal homology, and information loss}
\author{Technical research continuation}
\date{3 September 2026}

\begin{document}
\maketitle

\begin{center}
\fbox{\begin{minipage}{0.93\textwidth}\small
\textbf{Exact source basis.}
The source objects are the active Erd\H{o}s--Straus project archive of The
Clankers and the current canonical TeX/Bib bundle for the Alp\"oge--Fable/
Claude $S^6$ construction and its integrated audit.  The two supplied files
have SHA-256 values
\nolinkurl{d1e371380e9507a4e142dbf56430bc37145bf5ba42edfb76d21e898f0d131fea}
and
\nolinkurl{4a8453f2872dbd93269fc53c321cc19b74553fe0a4fa3f31c79644e28203c92a},
respectively.  Every map below displays its domain, codomain, inverse where it
exists, and lost data where it does not.  No regular $S^6$ torus fibre is
identified with a Koide point: the relevant coefficient matrix is singular.
\end{minipage}}
\end{center}

\begin{abstract}
For a prime $p\equiv1\pmod4$, every ordered Erd\H{o}s--Straus solution with
smallest denominator $a$ lies on a shell
\[
 Ryz=pa(y+z),\qquad 4a=p+R,\qquad R\equiv3\pmod4.
\]
The maintained divisor selector rewrites the shell by a divisor $u\mid a^2$
and one of two congruences.  We prove that the same selector state is exactly
an integral point on the incidence of two determinant loci:
\[
 \det\begin{pmatrix}R&-u\\4&-T\end{pmatrix}=-p^\varepsilon,
 \qquad
 \det\begin{pmatrix}u&a\\a&v\end{pmatrix}=0,
 \qquad uv=a^2.
\]
The first matrix is unimodularly row-equivalent to the terminal relation matrix
of the $(3,4,\infty)$ $S^6$ construction; the second is a positive integral
rank-one matrix and is exactly the Koide/$S^6$ determinant-boundary form.  A
unique factorization
\[
 u=hr^2,\qquad a=hrs,\qquad v=hs^2,\qquad \gcd(r,s)=1
\]
turns the two divisor channels into the linear congruences
$R\mid pr+s$ and $R\mid r+s$ and yields closed denominator formulas.  The pair
consisting of the rank-one matrix and terminal matrix is a globally injective
certificate; explicit inverse formulas recover $p,\varepsilon,j,u$ and the
unit-fraction denominators.  The exterior channel has terminal determinant
$-1$ and the printed integral homology package of $S^6$; the middle channel has
determinant $-p$ and $p$-torsion in degrees one through four.  The selector-to-solution map is also determined exactly: exterior states
are faithful, while middle states are paired by the fixed-point-free divisor
complement involution.  Finally, the full Erd\H{o}s--Straus conjecture is
equivalent to nonemptiness of the displayed two-determinant incidence for
every prime \(p\equiv1\pmod{12}\).  The construction does not prove that
nonemptiness.
\end{abstract}

\section{The fixed-prime shell}

Let $p\equiv1\pmod4$ be prime.  Write $\ES(p)$ for the assertion that there
exist positive integers $x\le y\le z$ with
\begin{equation}\label{eq:ES}
 \frac4p=\frac1x+\frac1y+\frac1z.
\end{equation}

\begin{proposition}[First-half denominator]\label{prop:first-half}
Every solution of \eqref{eq:ES} satisfies
\begin{equation}\label{eq:first-half}
 \frac p4<x\le\frac p2.
\end{equation}
Consequently there is a unique integer $j$ with
\begin{equation}\label{eq:R-a}
 0\le j<\frac{p-1}{4},\qquad
 R_j=4j+3,\qquad
 x=a_j=\frac{p+R_j}{4}.
\end{equation}
\end{proposition}

\begin{proof}
The strict lower bound follows from $1/x<4/p$.  Suppose $x>p/2$.
Clearing denominators gives
\[
 4xyz=p(xy+xz+yz),
\]
so $p\mid xyz$.  If $p\mid x$, all three denominators are at least $p$ and
the sum is at most $3/p$.  If $p\mid y$, then $y,z\ge p$ and
$1/x+1/y+1/z<2/p+2/p$.  Hence $z=kp$.  Substitution gives
\[
 (4k-1)xy=kp(x+y).
\]
Neither $x$ nor $y$ is divisible by $p$, so $4k-1=\ell p$.  Because
$p\equiv1\pmod4$, $\ell\equiv3\pmod4$ and $\ell\ge3$.  After cancelling $p$,
\[
 x\le\frac{2k}{\ell}
 =\frac p2+\frac1{2\ell}<\frac{p+1}{2},
\]
contradicting the integral inequality $x>p/2$.  Formula \eqref{eq:R-a} is the
unique enumeration of the integers in \eqref{eq:first-half}.
\end{proof}

Fix $j$ and abbreviate
\begin{equation}\label{eq:fixed-shell-data}
 R=R_j=4j+3,\qquad a=a_j=\frac{p+R}{4},\qquad N=pa.
\end{equation}
Then
\begin{equation}\label{eq:gcds}
 \gcd(R,p)=\gcd(R,a)=\gcd(p,a)=1.
\end{equation}
After fixing $x=a$, equation \eqref{eq:ES} is
\begin{equation}\label{eq:shell}
 Ryz=N(y+z),
\end{equation}
which is equivalent to
\begin{equation}\label{eq:factor-shell}
 (Ry-N)(Rz-N)=N^2.
\end{equation}
Thus, with
\begin{equation}\label{eq:d-e}
 d=Ry-N,
 \qquad e=Rz-N,
\end{equation}
we have $d,e>0$, $de=N^2$, and
\begin{equation}\label{eq:denom-inverse}
 y=\frac{N+d}{R},\qquad z=\frac{N+e}{R}.
\end{equation}
Conversely, any positive divisor $d\mid N^2$ satisfying
$d\equiv-N\pmod R$ gives \eqref{eq:shell} through
\eqref{eq:denom-inverse}; its complementary divisor automatically satisfies
the same congruence.

\section{The two complete divisor channels}

Because $\gcd(p,a)=1$, the $p$-adic exponent of $d$ is $0$, $1$, or $2$.
The involution $d\leftrightarrow e$ exchanges exponents $0$ and $2$.  Retain
the exponent-two representative in that pair and define
\begin{align}
 E_j&=\{u\in\Z_{>0}:u\mid a_j^2,\ R_j\mid4u+1\},\label{eq:E}\\
 M_j&=\{u\in\Z_{>0}:u\mid a_j^2,\ R_j\mid4u+p\}.\label{eq:M}
\end{align}
The letter $E$ records the exterior $p$-adic channel $d=p^2u$; the letter
$M$ records the middle channel $d=pu$.

\begin{theorem}[Complete two-channel selector]\label{thm:selector}
For $p\equiv1\pmod4$ prime,
\begin{equation}\label{eq:ES-selector}
 \ES(p)
 \quad\Longleftrightarrow\quad
 \bigsqcup_{0\le j<(p-1)/4}(E_j\sqcup M_j)\ne\varnothing.
\end{equation}
For a state $(\varepsilon,j,u)$, where $\varepsilon=0$ on $E_j$ and
$\varepsilon=1$ on $M_j$, put
\begin{equation}\label{eq:T-K}
 T=\frac{4u+p^\varepsilon}{R_j},
 \qquad
 K=\frac{p^{1-\varepsilon}T+1}{4}.
\end{equation}
Then $T,K\in\Z_{>0}$ and the corresponding denominators are
\begin{equation}\label{eq:source-denominators}
 (x,y,z)=
 \left(a_j,\frac{p^\varepsilon a_jK}{u},pK\right),
\end{equation}
up to interchange of $y,z$.
\end{theorem}

\begin{proof}
If $d=p^2u$, then $e=a^2/u$ and
\[
 R\mid d+N
 \Longleftrightarrow R\mid pu+a
 \Longleftrightarrow R\mid4u+1.
\]
If $d=pu$, then $e=pa^2/u$ and
\[
 R\mid d+N
 \Longleftrightarrow R\mid u+a
 \Longleftrightarrow R\mid4u+p.
\]
This proves completeness under the factor involution.  In either channel,
$RT=4u+p^\varepsilon$.  Since $R\equiv3\pmod4$ and
$p^\varepsilon\equiv1\pmod4$, one has $T\equiv3\pmod4$, so $K$ is integral.
The identities
\[
 RK=a+p^{1-\varepsilon}u,
 \qquad Ta=u+p^\varepsilon K
\]
give the divisibility in \eqref{eq:source-denominators}; direct substitution
then yields \eqref{eq:ES}.
\end{proof}

\section{The terminal determinant}

Define
\begin{equation}\label{eq:lambda}
 \lambda(\ell_0,\ell_1,\ell_2)
 =12\ell_0-4\ell_1-3\ell_2
\end{equation}
and the ordered terminal relation matrix
\begin{equation}\label{eq:Rterminal}
 \mathsf R(\ell_0,\ell_1,\ell_2)
 =\begin{pmatrix}3&-\ell_1\\4&\ell_2-4\ell_0\end{pmatrix},
 \qquad
 \det\mathsf R=-\lambda.
\end{equation}
For a selector state set
\begin{equation}\label{eq:ellmap}
 \ell_0=\frac{T+1}{4},
 \qquad
 \ell_1=u-jT,
 \qquad
 \ell_2=1.
\end{equation}

\begin{theorem}[Exact selector-to-terminal morphism]\label{thm:terminal}
For every prime $p\equiv1\pmod4$ and every selector state,
\begin{equation}\label{eq:lambda-pe}
 \lambda(\ell_0,\ell_1,1)=p^\varepsilon
\end{equation}
and
\begin{equation}\label{eq:row-equivalence}
 \begin{pmatrix}1&-j\\0&1\end{pmatrix}
 \begin{pmatrix}R_j&-u\\4&-T\end{pmatrix}
 =
 \begin{pmatrix}3&-\ell_1\\4&1-4\ell_0\end{pmatrix}
 =\mathsf R(\ell_0,\ell_1,1).
\end{equation}
The inverse is
\begin{equation}\label{eq:terminal-inverse}
 T=4\ell_0-1,
 \qquad
 u=\ell_1+jT.
\end{equation}
The crosswalk therefore needs only $p\equiv1\pmod4$, not the stronger
restriction $p\equiv1\pmod{12}$ used to isolate the hard prime class in the
source chapter.
\end{theorem}

\begin{proof}
Substituting \eqref{eq:ellmap}, $R_j=4j+3$, and
$R_jT=4u+p^\varepsilon$ gives
\[
 \lambda(\ell_0,\ell_1,1)
 =3(T+1)-4(u-jT)-3
 =(4j+3)T-4u=p^\varepsilon.
\]
Matrix multiplication proves \eqref{eq:row-equivalence};
\eqref{eq:terminal-inverse} is literal inversion.  The parity argument uses
only $p^\varepsilon\equiv1\pmod4$.  The $S^6$ finite-filling admissibility
conditions on this slice are also automatic: $\ell_2=1$ is odd, and
\[
 p^\varepsilon=\lambda(\ell_0,\ell_1,1)
 \equiv-\ell_1\pmod3.
\]
Because $p\ne3$, this forces $3\nmid\ell_1$ without assuming
$p\equiv1\pmod3$.
\end{proof}


Put
\begin{equation}\label{eq:selector-relation-matrix}
 \mathsf A_{p,\varepsilon,j,u}
 =\begin{pmatrix}R_j&-u\\4&-T\end{pmatrix}.
\end{equation}

\begin{proposition}[Exact adjugate factorization]
\label{prop:adjugate-factorization}
Every selector state satisfies
\begin{equation}\label{eq:adjugate-factorization}
 \mathsf A_{p,\varepsilon,j,u}
 \begin{pmatrix}T&-u\\4&-R_j\end{pmatrix}
 =p^\varepsilon I_2.
\end{equation}
Thus
\begin{equation}\label{eq:A-inverse}
 \mathsf A_{p,\varepsilon,j,u}^{-1}
 =\frac1{p^\varepsilon}
 \begin{pmatrix}T&-u\\4&-R_j\end{pmatrix}.
\end{equation}
For the primitive rank-one completion below,
\begin{equation}\label{eq:Gamma-primitive-kernel}
 \Gamma(a,u)\binom{s}{-r}=0,
 \qquad \gcd(r,s)=1.
\end{equation}
Hence a selector state carries simultaneously an index-\(p^\varepsilon\)
nonsingular relation matrix and a singular positive rank-one matrix with a
primitive integral kernel vector.
\end{proposition}

\begin{proof}
The four entries in the product on the left of
\eqref{eq:adjugate-factorization} are
\[
 R_jT-4u=p^\varepsilon,
 \qquad -R_ju+uR_j=0,
 \qquad 4T-4T=0,
 \qquad -4u+TR_j=p^\varepsilon.
\]
This proves the factorization and inverse.  Substitution of
\(\Gamma=h\binom r s(r\ s)\) gives
\(\Gamma(s,-r)^{\mathsf T}=0\), and primitivity is
\(\gcd(r,s)=1\).
\end{proof}

The left factor in \eqref{eq:row-equivalence} lies in
$\operatorname{SL}_2(\Z)$.  Hence both matrices present the same cokernel.
Since the terminal matrix contains entries $3$ and $4$, the gcd of all its
entries is one, and
\begin{equation}\label{eq:SNF}
 \SNF(\mathsf R)=\diag(1,p^\varepsilon),
 \qquad
 \coker\mathsf R\cong\Z/p^\varepsilon\Z.
\end{equation}
The numerical appearance of the pair $(3,4)$ is exact: the ES residue
$R_j=4j+3$ is changed to $3$ by subtracting $j$ times the second relation,
while the coefficient $4$ remains fixed.  These are the same two integers as
the finite orbifold multiplicities in the terminal $S^6$ presentation.

\section{The rank-one completion of the divisor gate}

Let $u\mid a^2$ and set
\begin{equation}\label{eq:v}
 v=\frac{a^2}{u}.
\end{equation}
Define
\begin{equation}\label{eq:Gamma}
 \Gamma(a,u)
 =\begin{pmatrix}u&a\\a&v\end{pmatrix}
 \in\Sym_2(\Z).
\end{equation}

\begin{proposition}[Divisibility as a unique rank-one completion]
\label{prop:completion}
For positive integers $a,u$, the following are equivalent:
\begin{enumerate}[label=(\roman*)]
\item $u\mid a^2$;
\item there is a positive integer $v$ for which
$\det\bigl(\begin{smallmatrix}u&a\\a&v\end{smallmatrix}\bigr)=0$.
\end{enumerate}
When these conditions hold, $v=a^2/u$ is unique, $\Gamma(a,u)$ is positive
semidefinite of rank one, and the map
\begin{equation}\label{eq:completion-map}
 (a,u)\longmapsto\Gamma(a,u)
\end{equation}
has inverse $\Gamma\mapsto(\Gamma_{12},\Gamma_{11})$ on its image.
\end{proposition}

\begin{proof}
The determinant equation is $uv-a^2=0$.  It has an integral positive solution
$v$ exactly when $u\mid a^2$, and then $v$ is forced.  The trace is positive
and the determinant is zero, so the eigenvalues are $u+v$ and zero.
\end{proof}

\begin{theorem}[Unique primitive-ray factorization]\label{thm:primitive}
Every matrix \eqref{eq:Gamma} has a unique factorization
\begin{equation}\label{eq:primitive-factorization}
 u=hr^2,
 \qquad
 a=hrs,
 \qquad
 v=hs^2,
 \qquad
 h,r,s\in\Z_{>0},
 \qquad
 \gcd(r,s)=1.
\end{equation}
Equivalently,
\begin{equation}\label{eq:Gamma-ray}
 \Gamma(a,u)=h\binom r s(r\ s),
 \qquad
 \SNF(\Gamma)=\diag(h,0).
\end{equation}
\end{theorem}

\begin{proof}
Put $g=\gcd(a,u)$, $r=u/g$, and $s=a/g$.  Then $\gcd(r,s)=1$.
The divisibility $u\mid a^2$ gives $gr\mid g^2s^2$, so
$r\mid gs^2$.  Coprimality gives $r\mid g$; put $h=g/r$.  Then
$u=hr^2$, $a=hrs$, and $v=a^2/u=hs^2$.  Conversely these equations clearly
produce \eqref{eq:Gamma}.  Any second primitive factorization has
$r/s=u/a$ in lowest terms, so it has the same $r,s$, and then the same $h$.
The gcd of all four matrix entries is $h$, while its determinant is zero,
which proves the Smith form.
\end{proof}

\section{The selector becomes two linear congruences}

\begin{theorem}[Primitive channel equations]\label{thm:primitive-channels}
Let \eqref{eq:primitive-factorization} hold.  Then
\begin{align}
 u\in E_j
 &\quad\Longleftrightarrow\quad R_j\mid pr+s,
 \label{eq:E-primitive}\\
 u\in M_j
 &\quad\Longleftrightarrow\quad R_j\mid r+s.
 \label{eq:M-primitive}
\end{align}
When \eqref{eq:E-primitive} holds, put
\begin{equation}\label{eq:kappaE}
 \kappa_E=\frac{pr+s}{R_j}.
\end{equation}
Then an ordered denominator triple is
\begin{equation}\label{eq:E-denominators}
 (x,y,z)=\bigl(hrs,\;phr\kappa_E,\;hs\kappa_E\bigr).
\end{equation}
When \eqref{eq:M-primitive} holds, put
\begin{equation}\label{eq:kappaM}
 \kappa_M=\frac{r+s}{R_j}.
\end{equation}
Then an ordered denominator triple is
\begin{equation}\label{eq:M-denominators}
 (x,y,z)=\bigl(hrs,\;phr\kappa_M,\;phs\kappa_M\bigr).
\end{equation}
Interchanging the last two entries gives the complementary orientation.
\end{theorem}

\begin{proof}
Modulo $R=R_j$, the identity $4a=p+R$ gives
$p\equiv4hrs$.  Since $R$ is coprime to $a=hrs$, each of $h,r,s$ is a unit
modulo $R$.  Thus
\[
 pr+s\equiv s(4hr^2+1)\pmod R,
\]
which proves \eqref{eq:E-primitive}.  Similarly
\[
 4u+p\equiv4hr^2+4hrs=4hr(r+s)\pmod R,
\]
which proves \eqref{eq:M-primitive}.

For \eqref{eq:E-denominators}, a common denominator calculation gives
\[
 \frac1{hrs}+\frac1{phr\kappa_E}+\frac1{hs\kappa_E}
 =\frac{p\kappa_E+s+pr}{phrs\kappa_E}
 =\frac{(p+R)\kappa_E}{phrs\kappa_E}
 =\frac4p.
\]
The middle-channel calculation is
\[
 \frac1{hrs}+\frac1{phr\kappa_M}+\frac1{phs\kappa_M}
 =\frac{p\kappa_M+s+r}{phrs\kappa_M}
 =\frac{(p+R)\kappa_M}{phrs\kappa_M}
 =\frac4p.
\]
\end{proof}

\begin{corollary}[Channel overlap]\label{cor:overlap}
Fix $(p,j,u)$ which lies in at least one of $E_j,M_j$.  It lies in both
channels exactly when
\begin{equation}\label{eq:overlap}
 R_j\mid p-1.
\end{equation}
\end{corollary}

\begin{proof}
The difference between the two primitive congruence numerators is
$(p-1)r$.  Since $r$ is a unit modulo $R_j$, the two congruences are
equivalent precisely under \eqref{eq:overlap}.  The hypothesis that one
already holds supplies the other.
\end{proof}


\section{Primitive-ray residues and the divisor-complement reflection}

Let
\begin{equation}\label{eq:swap-matrix}
 J=\begin{pmatrix}0&1\\1&0\end{pmatrix}.
\end{equation}
For a rank-one completion
\(
 \Gamma=h\binom r s(r\ s)
\)
write
\begin{equation}\label{eq:primitive-residue}
 \rho_{R}(\Gamma):=s r^{-1}\pmod R.
\end{equation}
This is defined because \(\gcd(R,hrs)=1\).

\begin{proposition}[The two channels as exact primitive-ray residues]
\label{prop:projective-rays}
For the fixed shell \((p,R,a)\),
\begin{align}
 u\in E_j
 &\quad\Longleftrightarrow\quad
 \rho_R(\Gamma)\equiv-p\pmod R,
 \label{eq:E-projective}\\
 u\in M_j
 &\quad\Longleftrightarrow\quad
 \rho_R(\Gamma)\equiv-1\pmod R.
 \label{eq:M-projective}
\end{align}
Thus the two channels are two exact residue points of the primitive rank-one
ray.  They coincide exactly when \(R\mid p-1\).
\end{proposition}

\begin{proof}
Theorem~\ref{thm:primitive-channels} gives
\(pr+s\equiv0\pmod R\) in the exterior channel and
\(r+s\equiv0\pmod R\) in the middle channel.  Multiplication by the inverse
of \(r\) proves \eqref{eq:E-projective} and \eqref{eq:M-projective}.  Their
right sides coincide exactly when \(p\equiv1\pmod R\).
\end{proof}

The divisor-complement map is
\begin{equation}\label{eq:divisor-complement}
 \sigma_j(u)=v=\frac{a_j^2}{u},
 \qquad
 \Gamma(a_j,v)=J\Gamma(a_j,u)J.
\end{equation}

\begin{theorem}[Exact complement action]
\label{thm:complement-action}
For every fixed shell:
\begin{enumerate}[label=(\roman*)]
\item \(\sigma_j\) preserves \(M_j\), has no fixed point there, and the two
states \(u\) and \(v\) reconstruct the same unordered denominator pair.
\item If \(u\in E_j\), then
\begin{equation}\label{eq:E-complement-condition}
 v\in E_j
 \quad\Longleftrightarrow\quad
 R_j\mid p^2-1.
\end{equation}
\item The diagonal rank-one matrix
\begin{equation}\label{eq:diagonal-gamma}
 \Gamma=a_j\begin{pmatrix}1&1\\1&1\end{pmatrix}
\end{equation}
occurs in the exterior channel exactly when
\begin{equation}\label{eq:diagonal-condition}
 R_j\mid p+1.
\end{equation}
In that case, with \(\kappa=(p+1)/R_j\), it gives
\begin{equation}\label{eq:diagonal-denominators}
 \frac4p=\frac1{a_j}+\frac1{pa_j\kappa}
             +\frac1{a_j\kappa}.
\end{equation}
\end{enumerate}
\end{theorem}

\begin{proof}
Under \(u=hr^2,a=hrs,v=hs^2\), the complement interchanges \(r\) and
\(s\).  The middle congruence \(R\mid r+s\) is symmetric, and
\eqref{eq:M-denominators} shows that the last two denominators are exchanged.
A fixed point would give \(u=v=a\), but then the middle condition would imply
\[
 R\mid4a+p=2p+R,
\]
hence \(R\mid2\), contradicting \(R\ge3\) odd and \(\gcd(R,p)=1\).

Suppose \(u\in E_j\).  Modulo \(R\),
\[
 4u\equiv-1,
 \qquad 4a\equiv p,
 \qquad
 4v=4a^2u^{-1}\equiv-p^2.
\]
Therefore \(4v+1\equiv1-p^2\pmod R\), proving
\eqref{eq:E-complement-condition}.  Finally \(u=v=a\) is an exterior state
exactly when
\[
 R\mid4a+1=p+R+1,
\]
which is \eqref{eq:diagonal-condition}; substitution of \(r=s=1,h=a\) into
\eqref{eq:E-denominators} gives \eqref{eq:diagonal-denominators}.
\end{proof}

Let \(\mathbf r_K(a,u)\) be the Koide vector defined below in
\eqref{eq:Koide-vector}.  The complement has an exact action already in the
original three Koide coordinates.

\begin{proposition}[Denominator exchange as a Koide reflection]
\label{prop:Koide-reflection}
Define
\begin{equation}\label{eq:WK}
 W_K=\frac13
 \begin{pmatrix}
  2&2&-1\\
  2&-1&2\\
  -1&2&2
 \end{pmatrix}.
\end{equation}
Then
\begin{equation}\label{eq:WK-properties}
 W_K^{\mathsf T}W_K=I,
 \qquad W_K^2=I,
 \qquad\det W_K=-1,
 \qquad Q_K(W_Kr)=Q_K(r),
\end{equation}
and
\begin{equation}\label{eq:WK-complement}
 \mathbf r_K(a,v)=W_K\mathbf r_K(a,u).
\end{equation}
In particular, middle-channel denominator exchange is represented exactly by
this rational orthogonal reflection on the Koide null cone.
\end{proposition}

\begin{proof}
Conjugation by \(J\) sends
\[
 \left(\frac{u+v}{2},a,\frac{u-v}{2}\right)
 \longmapsto
 \left(\frac{u+v}{2},a,-\frac{u-v}{2}\right).
\]
Thus in Lorentz coordinates the complement is
\(\operatorname{diag}(1,1,-1)\).  Conjugating by the orthogonal matrix
\(\mathcal F\) gives
\[
 W_K=\mathcal F^{-1}\operatorname{diag}(1,1,-1)\mathcal F,
\]
whose direct multiplication is \eqref{eq:WK}.  The identities in
\eqref{eq:WK-properties} and \eqref{eq:WK-complement} follow.
\end{proof}

\section{The exact terminal--Koide incidence set}

For $\varepsilon\in\{0,1\}$ define $\mathscr I_\varepsilon(p)$ to be the set
of tuples
\begin{equation}\label{eq:incidence-tuple}
 (j,\ell_0,\ell_1,\Gamma)
\end{equation}
with
\begin{equation}\label{eq:incidence-range}
 0\le j<\frac{p-1}{4},
 \qquad
 \Gamma=\begin{pmatrix}u&a\\a&v\end{pmatrix}
 \in\Sym_2(\Z),
 \qquad u,a,v>0,
\end{equation}
subject to
\begin{align}
 T&=4\ell_0-1>0,\label{eq:incidence-T}\\
 a&=\frac{p+4j+3}{4},\label{eq:incidence-a}\\
 u&=\ell_1+jT,\label{eq:incidence-u}\\
 \lambda(\ell_0,\ell_1,1)&=p^\varepsilon,\label{eq:incidence-lambda}\\
 \det\Gamma&=uv-a^2=0.\label{eq:incidence-zero}
\end{align}

\begin{theorem}[Erd\H{o}s--Straus as a two-determinant incidence problem]
\label{thm:incidence}
The map
\begin{equation}\label{eq:incidence-map}
 (\varepsilon,j,u)
 \longmapsto
 \left(j,\frac{T+1}{4},u-jT,
 \begin{pmatrix}u&a_j\\a_j&a_j^2/u\end{pmatrix}\right)
\end{equation}
is a bijection from the complete tagged selector to
$\mathscr I_0(p)\sqcup\mathscr I_1(p)$.  Its inverse is
\begin{equation}\label{eq:incidence-inverse}
 u=\Gamma_{11},
 \qquad
 \varepsilon=
 \begin{cases}
 0,&\lambda=1,\\
 1,&\lambda=p,
 \end{cases}
 \qquad
 j\text{ retained}.
\end{equation}
Consequently
\begin{equation}\label{eq:ES-incidence}
 \boxed{\ES(p)\Longleftrightarrow
 \mathscr I_0(p)\sqcup\mathscr I_1(p)\ne\varnothing.}
\end{equation}
Equations \eqref{eq:incidence-lambda} and \eqref{eq:incidence-zero} are
respectively the terminal determinant condition
\begin{equation}\label{eq:two-dets-1}
 \det\mathsf R(\ell_0,\ell_1,1)=-p^\varepsilon
\end{equation}
and the Koide/$S^6$ coefficient-boundary condition
\begin{equation}\label{eq:two-dets-2}
 \det\Gamma=0.
\end{equation}
\end{theorem}

\begin{proof}
The forward map is well defined by Theorems~\ref{thm:selector} and
\ref{thm:terminal} and Proposition~\ref{prop:completion}.  Conversely,
\eqref{eq:incidence-zero} gives $v=a^2/u\in\Z_{>0}$, hence $u\mid a^2$.
Expanding \eqref{eq:incidence-lambda} using
\eqref{eq:incidence-T}--\eqref{eq:incidence-u} gives
\[
 (4j+3)T-4u=p^\varepsilon,
\]
so $R_j\mid4u+p^\varepsilon$.  This is exactly the relevant selector
condition.  Both compositions are the identity entry by entry.
\end{proof}

The theorem identifies the surviving arithmetic difficulty exactly.  The
terminal determinant equation by itself has many integral solutions.  The
Erd\H{o}s--Straus gate is the additional requirement that the shared top row
$(u,a)$ admit the positive integral rank-one completion \eqref{eq:Gamma}.


\section{Exact quotient from selector states to denominator triples}

Let \(\mathscr S(p)\) be the set of triples
\((x,y,z)\in\Z_{>0}^3\) satisfying \eqref{eq:ES} and
\(x\le y\le z\).  Put
\begin{equation}\label{eq:EM-global-sets}
 \mathcal E(p)=\bigsqcup_j\{j\}\times E_j,
 \qquad
 \mathcal M(p)=\bigsqcup_j\{j\}\times M_j.
\end{equation}
By Theorem~\ref{thm:complement-action},
\begin{equation}\label{eq:M-involution}
 \sigma(j,u)=\left(j,\frac{a_j^2}{u}\right)
\end{equation}
is a fixed-point-free involution of \(\mathcal M(p)\).

\begin{theorem}[Exact solution-set quotient]
\label{thm:solution-quotient}
Sorting the three denominators in
\eqref{eq:E-denominators} and \eqref{eq:M-denominators} induces a bijection
\begin{equation}\label{eq:solution-quotient}
 \boxed{
 \mathcal E(p)\;\sqcup\;
 \bigl(\mathcal M(p)/\langle\sigma\rangle\bigr)
 \xrightarrow{\ \sim\ }
 \mathscr S(p).}
\end{equation}
Consequently
\begin{equation}\label{eq:solution-count}
 |\mathscr S(p)|
 =\sum_j|E_j|+\frac12\sum_j|M_j|.
\end{equation}
\end{theorem}

\begin{proof}
Surjectivity is Theorem~\ref{thm:selector} together with the unique
first-denominator index of Proposition~\ref{prop:first-half}.  Given a sorted
solution, the value \(x=a_j\) recovers \(j\).  The unordered residual factors
\[
 \{d,e\}=\{R_jy-pa_j,R_jz-pa_j\}
\]
have \(p\)-adic exponent multiset \(\{2,0\}\) in the exterior channel and
\(\{1,1\}\) in the middle channel; hence the channel is recovered.

In the exterior channel, the factor of \(p\)-adic exponent two uniquely has
form \(p^2u\), so it recovers \(u\).  Two different exterior states therefore
cannot give the same sorted triple.  In the middle channel the unordered
factors are \(\{pu,pv\}\), so the sorted triple recovers exactly the orbit
\(\{u,v\}\).  The involution has no fixed point by
Theorem~\ref{thm:complement-action}.  This proves both injectivity and
\eqref{eq:solution-count}.
\end{proof}

\section{A faithful double certificate}

Define the global tagged selector set by taking the disjoint union over all
primes $p\equiv1\pmod4$.  The double-certificate map is
\begin{equation}\label{eq:double-cert}
 \mathfrak C(p,\varepsilon,j,u)
 =\left(
 \Gamma(a_j,u),
 \mathsf R(\ell_0,\ell_1,1)
 \right).
\end{equation}

\begin{theorem}[Global injectivity and explicit inverse]\label{thm:faithful}
The map \eqref{eq:double-cert} is injective.  For a pair in its image, write
\[
 \Gamma=\begin{pmatrix}u&a\\a&v\end{pmatrix},
 \qquad
 \mathsf R=\begin{pmatrix}3&-\ell_1\\4&-T\end{pmatrix}.
\]
Then the inverse is
\begin{equation}\label{eq:faithful-inverse}
 j=\frac{u-\ell_1}{T},
 \qquad
 p=4a-(4j+3),
 \qquad
 \lambda=-\det\mathsf R,
 \qquad
 \varepsilon=
 \begin{cases}0,&\lambda=1,\\1,&\lambda=p.\end{cases}
\end{equation}
The complementary divisor is $v=\Gamma_{22}$, and
Theorem~\ref{thm:primitive-channels} then recovers the denominator triple.
\end{theorem}

\begin{proof}
Every formula in \eqref{eq:faithful-inverse} follows from a displayed entry of
the two matrices: $u=\Gamma_{11}$, $a=\Gamma_{12}$,
$T=-\mathsf R_{22}$, and $\ell_1=-\mathsf R_{12}$.  The source equation
$\ell_1=u-jT$ gives $j$; $4a=p+4j+3$ gives $p$; the determinant gives the
channel.  Thus two selector states with the same matrix pair have identical
$p,\varepsilon,j,u$.
\end{proof}

The terminal matrix alone is not faithful.  At $p=13$, the exterior states
$(j,u)=(0,2)$ and $(1,5)$ both give
\begin{equation}\label{eq:same-terminal}
 \mathsf R=\begin{pmatrix}3&-2\\4&-3\end{pmatrix},
\end{equation}
but their rank-one matrices are respectively
\begin{equation}\label{eq:different-gamma}
 \begin{pmatrix}2&4\\4&8\end{pmatrix},
 \qquad
 \begin{pmatrix}5&5\\5&5\end{pmatrix}.
\end{equation}
The same terminal matrix also occurs at $p=73$, $j=1$, $u=5$, where
\begin{equation}\label{eq:p73-gamma}
 \Gamma=\begin{pmatrix}5&20\\20&80\end{pmatrix}.
\end{equation}
The matrix pair restores exactly the information lost by the terminal
projection.

\section{The Koide and \texorpdfstring{$S^6$}{S6} coefficient maps}

Define the native $S^6$ coefficient form
\begin{equation}\label{eq:Gnative}
 G(t,m,b)=\begin{pmatrix}2t&2m\\2m&b/3\end{pmatrix}
\end{equation}
and the native Lorentz dictionary
\begin{equation}\label{eq:Lnative}
 \mathcal L(t,m,b)
 =\left(t+\frac b6,\;2m,\;t-\frac b6\right)=(H,X,Y).
\end{equation}
Its inverse is
\begin{equation}\label{eq:Linv}
 \mathcal L^{-1}(H,X,Y)
 =\left(\frac{H+Y}{2},\frac X2,3(H-Y)\right).
\end{equation}
The raw-period dictionary is
\begin{equation}\label{eq:PK}
 \mathcal P_K(H,X,Y)
 =\left(H+Y,\frac{X}{\sqrt6},H-Y\right)
 =(T_{\rm raw},M_{\rm raw},B_{\rm raw}).
\end{equation}
The exact morphism between the two coefficient triples is
\begin{equation}\label{eq:Dmap}
 \mathcal D(t,m,b)
 =\left(2t,\sqrt{\frac23}\,m,\frac b3\right),
 \qquad
 \mathcal D^{-1}(T,M,B)
 =\left(\frac T2,\sqrt{\frac32}\,M,3B\right),
\end{equation}
and
\begin{equation}\label{eq:dictionary-commutes}
 \mathcal P_K=\mathcal D\circ\mathcal L^{-1}.
\end{equation}

\begin{theorem}[Selector state to the Koide--$S^6$ boundary]
\label{thm:koide-boundary}
For a selector state, put $v=a^2/u$ and define
\begin{equation}\label{eq:S6-coeff-state}
 (t,m,b)=\left(\frac u2,\frac a2,3v\right).
\end{equation}
Then
\begin{equation}\label{eq:G=Gamma}
 G(t,m,b)=\Gamma(a,u),
\end{equation}
\begin{equation}\label{eq:native-Lorentz-state}
 \mathcal L(t,m,b)
 =\left(\frac{u+v}{2},a,\frac{u-v}{2}\right),
\end{equation}
and
\begin{equation}\label{eq:raw-state}
 \mathcal D(t,m,b)=\left(u,\frac a{\sqrt6},v\right).
\end{equation}
Therefore
\begin{align}
 \det G(t,m,b)
 &=\frac23(tb-6m^2)=0,\label{eq:native-zero}\\
 T_{\rm raw}B_{\rm raw}-6M_{\rm raw}^2
 &=uv-a^2=0.\label{eq:raw-zero}
\end{align}
\end{theorem}

\begin{proof}
Each assertion is direct substitution.  In particular,
\[
 G\left(\frac u2,\frac a2,3v\right)
 =\begin{pmatrix}u&a\\a&v\end{pmatrix}.
\]
Equation \eqref{eq:dictionary-commutes} follows by applying
\eqref{eq:Dmap} to \eqref{eq:Linv}.
\end{proof}

Let
\begin{equation}\label{eq:QK}
 Q_K(r_1,r_2,r_3)
 =\frac23(r_1+r_2+r_3)^2-(r_1^2+r_2^2+r_3^2)
\end{equation}
and let
\begin{equation}\label{eq:F}
 \mathcal F(r_1,r_2,r_3)
 =\left(
 \frac{r_1+r_2+r_3}{\sqrt3},
 \frac{r_1-r_3}{\sqrt2},
 \frac{r_1+r_3-2r_2}{\sqrt6}
 \right)=(H,X,Y).
\end{equation}
Then $\mathcal F$ is orthogonal and
\begin{equation}\label{eq:QK-Lorentz}
 Q_K=H^2-X^2-Y^2.
\end{equation}
The exact Koide vector associated with the selector matrix is
\begin{equation}\label{eq:Koide-vector}
 \mathbf r_K(a,u)=
 \mathcal F^{-1}\left(\frac{u+v}{2},a,\frac{u-v}{2}\right),
\end{equation}
that is,
\begin{align}
 (\mathbf r_K)_1
 &=\frac{u+v}{2\sqrt3}+\frac a{\sqrt2}
   +\frac{u-v}{2\sqrt6},\label{eq:rK1}\\
 (\mathbf r_K)_2
 &=\frac{u+v}{2\sqrt3}-\frac{u-v}{\sqrt6},\label{eq:rK2}\\
 (\mathbf r_K)_3
 &=\frac{u+v}{2\sqrt3}-\frac a{\sqrt2}
   +\frac{u-v}{2\sqrt6}.
 \label{eq:rK3}
\end{align}
It satisfies
\begin{equation}\label{eq:rKnull}
 Q_K(\mathbf r_K(a,u))=0.
\end{equation}

\begin{proposition}[The arithmetic field is necessarily enlarged]
\label{prop:anisotropic}
The equation $Q_K(r)=0$ has no nonzero solution in $\Q^3$.  Thus every
nonzero integral selector matrix gives a Lorentz-rational null point through
\eqref{eq:native-Lorentz-state}, but its Koide coordinates
\eqref{eq:rK1}--\eqref{eq:rK3} lie in $\Kfield^3$ and never in
$\Q^3\setminus\{0\}$.
\end{proposition}

\begin{proof}
For rational $r_i$, clear denominators and put
\[
 S=r_1+r_2+r_3,
 \qquad U=r_1-r_2,
 \qquad V=r_1+r_2-2r_3.
\]
A direct invertible rational change of coordinates gives
\begin{equation}\label{eq:rational-Q}
 Q_K=\frac{2S^2-3U^2-V^2}{6}.
\end{equation}
Suppose a primitive nonzero integer triple satisfies
$2S^2=3U^2+V^2$.  Modulo $3$, the equality
$V^2\equiv2S^2$ forces $3\mid S,V$.  Dividing by $3$ and reducing again
modulo $3$ forces $3\mid U$.  This contradicts primitivity.  The field
claim follows from the explicit formulas.
\end{proof}

\begin{proposition}[Exact positive Koide sector]\label{prop:positive-sector}
Put
\begin{equation}\label{eq:theta}
 \theta=\frac ua=\frac rs>0
\end{equation}
and
\begin{equation}\label{eq:thresholds}
 \theta_-=
 \frac{\sqrt3-\sqrt2}{1+\sqrt2},
 \qquad
 \theta_+=\frac{\sqrt3+\sqrt2}{1+\sqrt2}.
\end{equation}
The first coordinate in \eqref{eq:rK1} is always positive, and all three
Koide coordinates are positive exactly when
\begin{equation}\label{eq:positive-sector-condition}
 0<\theta<\theta_-
 \qquad\text{or}\qquad
 \theta_+<\theta<1+\sqrt2.
\end{equation}
\end{proposition}

\begin{proof}
Use $u=a\theta$ and $v=a/\theta$ in
\eqref{eq:rK1}--\eqref{eq:rK3}.  After multiplication by the positive factor
$2\sqrt6\theta/a$, the second coordinate has sign
\[
 (\sqrt2-2)\theta^2+(\sqrt2+2),
\]
which is positive exactly for $\theta<1+\sqrt2$.  The third has sign
\[
 (1+\sqrt2)\theta^2-2\sqrt3\theta+(\sqrt2-1),
\]
whose roots are $\theta_-,\theta_+$ and whose leading coefficient is
positive.  The corresponding expression for the first coordinate has all
positive coefficients.  Because $\theta\in\Q$ and the three boundary numbers
are irrational, equality does not occur.
\end{proof}

\begin{remark}[Exact image obstruction]
The map \eqref{eq:S6-coeff-state} lands in the determinant-zero boundary of
the ambient coefficient space.  For every choice of the real parts of
$\tau,\mu,\beta$, the canonical real period matrix with these imaginary
coefficients satisfies $\det P=0$.  Its four columns therefore do not generate
a full lattice in $\R^4$.  The map is not a map to a regular complex-torus
fibre.  The actual radial $S^6$ cusp instead has
$\det P_R=-\delta_R\ne0$ with $\delta_R\sim R^2$.
\end{remark}

\section{The full residual shell as a second rank-one matrix}

The selector matrix $\Gamma$ records $u,a,v$.  The complete two-denominator
shell records $d,N,e$ in
\begin{equation}\label{eq:shell-matrix}
 \Gamma_{\rm sh}=\begin{pmatrix}d&N\\N&e\end{pmatrix},
 \qquad de=N^2.
\end{equation}
The exact map from the selector matrix is channel dependent:
\begin{align}
 \Gamma_{\rm sh}^{E}
 &=\begin{pmatrix}p&0\\0&1\end{pmatrix}
   \Gamma
   \begin{pmatrix}p&0\\0&1\end{pmatrix}
 =\begin{pmatrix}p^2u&pa\\pa&v\end{pmatrix},
 \label{eq:shell-E}\\
 \Gamma_{\rm sh}^{M}
 &=p\Gamma
 =\begin{pmatrix}pu&pa\\pa&pv\end{pmatrix}.
 \label{eq:shell-M}
\end{align}

Put
\begin{equation}\label{eq:selector-lorentz-vector}
 \mathbf L(\Gamma)=
 \begin{pmatrix}H\\X\\Y\end{pmatrix}
 =\begin{pmatrix}(u+v)/2\\a\\(u-v)/2\end{pmatrix},
 \qquad
 Q_L(H,X,Y)=H^2-X^2-Y^2.
\end{equation}

\begin{theorem}[The two ES channels as exact Lorentz-conformal maps]
\label{thm:channel-lorentz-maps}
The channel maps in \eqref{eq:shell-E}--\eqref{eq:shell-M} induce the following
linear maps on the Lorentz coordinates of the rank-one symmetric matrix:
\begin{align}
 \mathbf L(\Gamma_{\rm sh}^{E})
 &=\mathcal B_p\mathbf L(\Gamma),
 &\mathcal B_p
 &=\begin{pmatrix}
 \dfrac{p^2+1}{2}&0&\dfrac{p^2-1}{2}\\
 0&p&0\\
 \dfrac{p^2-1}{2}&0&\dfrac{p^2+1}{2}
 \end{pmatrix},
 \label{eq:E-Lorentz-map}\\
 \mathbf L(\Gamma_{\rm sh}^{M})
 &=pI_3\mathbf L(\Gamma).
 \label{eq:M-Lorentz-map}
\end{align}
With \(J_L=\diag(1,-1,-1)\),
\begin{equation}\label{eq:conformal-Lorentz}
 \mathcal B_p^{\mathsf T}J_L\mathcal B_p=p^2J_L,
 \qquad
 (pI_3)^{\mathsf T}J_L(pI_3)=p^2J_L.
\end{equation}
Consequently
\begin{equation}\label{eq:channel-Q-scaling}
 Q_L(\mathbf L(\Gamma_{\rm sh}^{E}))
 =Q_L(\mathbf L(\Gamma_{\rm sh}^{M}))
 =p^2Q_L(\mathbf L(\Gamma)).
\end{equation}
On primitive rank-one vectors, the exterior channel sends
\begin{equation}\label{eq:E-ray-map}
 h\binom r s(r\ s)
 \longmapsto
 h\binom{pr}{s}(pr\ s),
\end{equation}
whereas the middle channel sends
\begin{equation}\label{eq:M-ray-map}
 h\binom r s(r\ s)
 \longmapsto
 ph\binom r s(r\ s).
\end{equation}
Thus the exterior channel changes the primitive projective ray
\([r:s]\) to \([pr:s]\), while the middle channel retains \([r:s]\) and
multiplies its content by \(p\).
\end{theorem}

\begin{proof}
Under the symmetric-matrix identification
\[
 \begin{pmatrix}H+Y&X\\X&H-Y\end{pmatrix},
\]
conjugation by \(\diag(p,1)\) gives
\[
 \begin{pmatrix}p^2(H+Y)&pX\\pX&H-Y\end{pmatrix}.
\]
Taking its half-trace, off-diagonal entry, and half diagonal difference gives
\eqref{eq:E-Lorentz-map}.  Multiplication of the symmetric matrix by \(p\)
gives \eqref{eq:M-Lorentz-map}.  Direct matrix multiplication proves
\eqref{eq:conformal-Lorentz}; the quadratic identity and the two primitive-ray
formulas follow.
\end{proof}

\begin{corollary}[Exact boost--dilation factorization]
\label{cor:boost-dilation}
Let
\begin{equation}\label{eq:boost-generator}
 \mathcal K_{HY}=
 \begin{pmatrix}0&0&1\\0&0&0\\1&0&0\end{pmatrix}.
\end{equation}
Then
\begin{equation}\label{eq:boost-generator-lorentz}
 \mathcal K_{HY}^{\mathsf T}J_L+J_L\mathcal K_{HY}=0
\end{equation}
and the exterior channel matrix has the exact factorization
\begin{equation}\label{eq:boost-factorization}
 \mathcal B_p
 =p\exp\!\bigl((\log p)\mathcal K_{HY}\bigr)
 =p\begin{pmatrix}
 \cosh(\log p)&0&\sinh(\log p)\\
 0&1&0\\
 \sinh(\log p)&0&\cosh(\log p)
 \end{pmatrix}.
\end{equation}
The middle channel is the pure dilation $pI_3$.  If
$\theta=s/r$ is the positive primitive slope, then the exterior channel sends
\begin{equation}\label{eq:slope-exterior}
 \theta\longmapsto\frac{\theta}{p},
\end{equation}
while the middle channel retains $\theta$.  With
\begin{equation}\label{eq:zeta-middle-function}
 \zeta_M(\theta)=-\frac{\theta}{(1+\theta)^2},
\end{equation}
one has the exact identity
\begin{equation}\label{eq:zeta-channel-composition}
 \zeta_E(\theta)=\zeta_M\!\left(\frac{\theta}{p}\right).
\end{equation}
\end{corollary}

\begin{proof}
Since
\[
 \cosh(\log p)=\frac{p+p^{-1}}2,
 \qquad
 \sinh(\log p)=\frac{p-p^{-1}}2,
\]
multiplication by $p$ gives the three entries of $\mathcal B_p$ in
\eqref{eq:E-Lorentz-map}.  Direct multiplication proves
\eqref{eq:boost-generator-lorentz}.  The primitive-ray map
$[r:s]\mapsto[pr:s]$ gives \eqref{eq:slope-exterior}, and
\[
 -\frac{p\theta}{(p+\theta)^2}
 =-\frac{\theta/p}{(1+\theta/p)^2}
\]
proves \eqref{eq:zeta-channel-composition}.
\end{proof}

Both matrices have determinant zero.  Their Sarnak coordinates are
\begin{equation}\label{eq:Sarnak}
 T_S=\frac{d+e}{2},
 \qquad X_S=N,
 \qquad Y_S=0,
 \qquad W_S=\frac{d-e}{2},
\end{equation}
and satisfy
\begin{equation}\label{eq:Sarnak-null}
 X_S^2+W_S^2-T_S^2=N^2-de=0.
\end{equation}
Denominator exchange is $W_S\mapsto-W_S$.

Define the exact quarter coordinate
\begin{equation}\label{eq:zeta}
 \zeta=-\frac{yz}{(y+z)^2}.
\end{equation}
Then
\begin{equation}\label{eq:quarter-general}
 \zeta=-\frac14+\frac14
 \left(\frac{W_S}{T_S+X_S}\right)^2.
\end{equation}
Using Theorem~\ref{thm:primitive-channels}, this becomes
\begin{align}
 \zeta_E&=-\frac{prs}{(pr+s)^2},\label{eq:zetaE}\\
 \zeta_M&=-\frac{rs}{(r+s)^2}.\label{eq:zetaM}
\end{align}

\begin{proposition}[The balanced quarter point is excluded]
\label{prop:quarter-excluded}
No Erd\H{o}s--Straus shell in these channels has $\zeta=-1/4$.
\end{proposition}

\begin{proof}
Equality in \eqref{eq:quarter-general} is equivalent to $d=e$, or $y=z$.
In the exterior channel this gives $pr=s$.  Primitivity forces $r=1,s=p$,
but then $R\mid pr+s=2p$, impossible because $R\ge3$ is odd and
$\gcd(R,p)=1$.  In the middle channel equality gives $r=s$, so primitivity
forces $r=s=1$, but then $R\mid r+s=2$, again impossible.
\end{proof}

The null cone \eqref{eq:Sarnak-null} is the $p=1$ member of the
Kocik--Sarnak four-circle family used in the Erd\H{o}s--Straus archive.  The
Kocik charged-lepton specialization is the different four-circle member
$p=2/3$, whose equation has coefficient $5/3$.  The rigorous Koide relation in
this note is instead the explicit three-dimensional Lorentz isometry
\eqref{eq:F}--\eqref{eq:QK-Lorentz}.  These two statements are connected by
the displayed maps, not by identifying $1$ with $2/3$.

\section{The terminal topological dichotomy}

The source terminal group is
\begin{equation}\label{eq:terminal-group}
 \mathcal G(\ell)=
 \langle c,x,y\mid c\text{ central},\ xy=c^{\ell_0},\
 x^3=c^{\ell_1},\ y^4=c^{\ell_2}\rangle.
\end{equation}
Eliminating $y$ makes the presentation abelian with relation matrix
\eqref{eq:Rterminal}.  Therefore
\begin{equation}\label{eq:group-outcome}
 \mathcal G(\ell)\cong\Z/p^\varepsilon\Z.
\end{equation}
The source all-degree terminal package gives, for $\lambda=p^\varepsilon$,
\begin{equation}\label{eq:homology-package}
 (H_0,\ldots,H_6)
 \cong
 \left(
 \Z,\Z/p^\varepsilon,\Z/p^\varepsilon,
 \Z/p^\varepsilon,\Z/p^\varepsilon,0,\Z
 \right).
\end{equation}
Here $\Z/1=0$.  Consequently:
\begin{center}
\begin{tabular}{c c c c}
\toprule
Channel & $\lambda$ & terminal group & terminal homology package\\
\midrule
$E$ & $1$ & $0$ & $H_*(S^6;\Z)$\\
$M$ & $p$ & $\Z/p$ & $\Z/p$ in degrees $1,2,3,4$\\
\bottomrule
\end{tabular}
\end{center}

Thus an $E$-channel Erd\H{o}s--Straus witness lands on the homology-sphere
terminal slice of the $S^6$ algebra.  An $M$-channel witness lands on a
$p$-torsion terminal slice.  A state in both channels supplies two different
terminal matrices over the same rank-one Koide boundary matrix.

\section{Worked exact examples}

\subsection{\texorpdfstring{$p=5$: extension beyond $p\equiv1\pmod{12}$}{p=5: extension beyond p=1 mod 12}}
Here $j=0$, $R=3$, $a=2$, and $u=2$ lies in $E_0$.  One has
\[
 \Gamma=\begin{pmatrix}2&2\\2&2\end{pmatrix},
 \qquad
 T=3,
 \qquad
 (\ell_0,\ell_1,\ell_2)=(1,2,1),
\]
\[
 \mathsf R=\begin{pmatrix}3&-2\\4&-3\end{pmatrix},
 \qquad
 \det\mathsf R=-1,
\]
and
\[
 \frac45=\frac12+\frac1{20}+\frac14.
\]
This proves by an explicit state that the terminal crosswalk itself applies to
the residue class $p\equiv5\pmod{12}$.

\subsection{\texorpdfstring{$p=13$: one rank-one ray, two channels}{p=13: one rank-one ray, two channels}}
For $j=0$, $R=3$, $a=4$, $u=2$, one has
\[
 \Gamma=\begin{pmatrix}2&4\\4&8\end{pmatrix}
 =2\binom12(1\ 2).
\]
The exterior state gives
\[
 T_E=3,
 \quad
 \mathsf R_E=\begin{pmatrix}3&-2\\4&-3\end{pmatrix},
 \quad
 \frac4{13}=\frac14+\frac1{130}+\frac1{20}.
\]
The same $(p,j,u)$ also satisfies the middle congruence because $3\mid12$:
\[
 T_M=7,
 \quad
 \mathsf R_M=\begin{pmatrix}3&-2\\4&-7\end{pmatrix},
 \quad
 \frac4{13}=\frac14+\frac1{26}+\frac1{52}.
\]
Thus the same Koide/$S^6$ boundary ray supports both the sphere-type terminal
package and the $13$-torsion terminal package.

For $j=1$, $R=7$, $a=5$, $u=5$, the exterior state has
\[
 \Gamma=\begin{pmatrix}5&5\\5&5\end{pmatrix},
 \qquad
 \mathsf R=\begin{pmatrix}3&-2\\4&-3\end{pmatrix},
 \qquad
 \frac4{13}=\frac15+\frac1{130}+\frac1{10}.
\]
This is the same terminal matrix as the preceding exterior state but a
different rank-one matrix; Theorem~\ref{thm:faithful} distinguishes them.


\section{The full Erd\H{o}s--Straus conjecture}

Let \(\mathrm{ES}\) denote the assertion that
\begin{equation}\label{eq:global-ES}
 \frac4n=\frac1x+\frac1y+\frac1z
\end{equation}
has a positive integral solution for every integer \(n\ge2\).

\begin{theorem}[Global reduction to the two-determinant incidence]
\label{thm:global-incidence}
The full conjecture is equivalent to
\begin{equation}\label{eq:global-incidence}
 \boxed{
 \mathrm{ES}
 \quad\Longleftrightarrow\quad
 \forall p\in\mathbb P,
 \ p\equiv1\pmod{12}:
 \ \mathscr I_0(p)\sqcup\mathscr I_1(p)\ne\varnothing.}
\end{equation}
Equivalently,
\begin{equation}\label{eq:global-counterexample}
 \neg\mathrm{ES}
 \quad\Longleftrightarrow\quad
 \exists p\in\mathbb P,
 \ p\equiv1\pmod{12}:
 \ \mathscr I_0(p)=\mathscr I_1(p)=\varnothing.
\end{equation}
The sphere-type subproblem
\begin{equation}\label{eq:E-only-sufficient}
 \forall p\equiv1\pmod{12},\quad
 \mathscr I_0(p)\ne\varnothing
\end{equation}
is sufficient for \(\mathrm{ES}\); it is not asserted to be necessary.
\end{theorem}

\begin{proof}
If \eqref{eq:global-ES} holds for a divisor \(d\mid n\), multiplication of all
three denominators by \(n/d\) gives a solution for \(n\).  It therefore
suffices, and is necessary, to treat primes.

For \(p=2\),
\[
 \frac42=\frac11+\frac12+\frac12.
\]
For a prime \(p\equiv3\pmod4\), put
\[
 a=\frac{p+1}{4},\qquad N=pa,\qquad
 b=N+1,\qquad c=N(N+1).
\]
Then
\[
 \frac4p-\frac1a=\frac1N
 =\frac1{N+1}+\frac1{N(N+1)}.
\]
For a prime \(p\equiv5\pmod{12}\), put
\[
 a=\frac{p+3}{4},\qquad
 b=\frac{(p+1)(p+3)}{12},\qquad c=pb.
\]
Then
\[
 \frac4p-\frac1a=\frac3{pa}
 =\frac1b+\frac1{pb}.
\]
Every prime greater than three is congruent to \(1,5,7\), or \(11\) modulo
\(12\); the last two classes are \(3\pmod4\).  Hence only primes
\(p\equiv1\pmod{12}\) remain.  For each of those primes,
Theorem~\ref{thm:incidence} is exactly the asserted incidence criterion.
Taking the logical negation gives \eqref{eq:global-counterexample}, and
retaining only the \(\varepsilon=0\) alternative gives
\eqref{eq:E-only-sufficient}.
\end{proof}

\begin{corollary}[Global primitive Koide-ray formulation]
\label{cor:global-primitive}
The full Erd\H{o}s--Straus conjecture is equivalent to the following statement:
for every prime \(p\equiv1\pmod{12}\), there exist
\(h,r,s\in\Z_{>0}\) such that
\begin{equation}\label{eq:global-primitive-data}
 \gcd(r,s)=1,
 \qquad
 \frac p4<hrs<\frac p2,
 \qquad
 R:=4hrs-p,
\end{equation}
and at least one of the two exact ray conditions holds:
\begin{equation}\label{eq:global-primitive-rays}
 R\mid pr+s
 \qquad\text{or}\qquad
 R\mid r+s.
\end{equation}
A hard prime \(p\equiv1\pmod{12}\) is a counterexample exactly when every
triple satisfying \eqref{eq:global-primitive-data} satisfies
\begin{equation}\label{eq:global-primitive-obstruction}
 R\nmid pr+s
 \qquad\text{and}\qquad
 R\nmid r+s.
\end{equation}
\end{corollary}

\begin{proof}
Given an incidence point, Theorem~\ref{thm:primitive} writes
\(u=hr^2,a=hrs,v=hs^2\).  Proposition~\ref{prop:first-half} gives the strict
interval in \eqref{eq:global-primitive-data}, and
Theorem~\ref{thm:primitive-channels} gives
\eqref{eq:global-primitive-rays}.  Conversely, a triple satisfying
\eqref{eq:global-primitive-data} has
\[
 0<R<p,
 \qquad R\equiv3\pmod4,
 \qquad j=\frac{R-3}{4}\in\Z,
 \qquad a=hrs=\frac{p+R}{4}.
\]
Set \(u=hr^2\) and \(v=hs^2\).  Then \(uv=a^2\), and either divisibility in
\eqref{eq:global-primitive-rays} is, by
Theorem~\ref{thm:primitive-channels}, exactly one of the two selector channels.
The global equivalence follows from Theorem~\ref{thm:global-incidence}; its
logical negation is \eqref{eq:global-primitive-obstruction}.
\end{proof}

\begin{corollary}[Divisor-free hard-prime formulation]
\label{cor:global-divisor-free}
The full Erd\H{o}s--Straus conjecture is equivalent to the following statement:
for every prime $p\equiv1\pmod{12}$ there exist positive coprime integers
$r,s$ and an integer $R$ such that
\begin{equation}\label{eq:divisor-free-range}
 0<R<p,
 \qquad R\equiv3\pmod4,
 \qquad 4rs\mid p+R,
\end{equation}
and
\begin{equation}\label{eq:divisor-free-rays}
 R\mid pr+s
 \qquad\text{or}\qquad
 R\mid r+s.
\end{equation}
For such data the content and first denominator are uniquely
\begin{equation}\label{eq:divisor-free-h-a}
 h=\frac{p+R}{4rs},
 \qquad a=hrs=\frac{p+R}{4}.
\end{equation}
\end{corollary}

\begin{proof}
The data of Corollary~\ref{cor:global-primitive} give
$R=4hrs-p$, hence $p+R=4hrs$ and therefore
$4rs\mid p+R$.  Conversely, \eqref{eq:divisor-free-range} defines the positive
integer $h$ in \eqref{eq:divisor-free-h-a}; then
$R=4hrs-p$ and $p/4<hrs<p/2$.  The two divisibility conditions are unchanged,
so Corollary~\ref{cor:global-primitive} applies in both directions.
\end{proof}

\begin{proposition}[Exact middle-channel shell bound]
\label{prop:middle-shell-bound}
Suppose a hard-prime selector lies in the middle channel and write
\begin{equation}\label{eq:middle-kappa}
 r+s=\kappa R,
 \qquad \kappa\in\Z_{>0}.
\end{equation}
Then
\begin{equation}\label{eq:middle-shell-bound}
 (4\kappa-1)R\le p+4.
\end{equation}
In particular
\begin{equation}\label{eq:middle-shell-one-third}
 R\le\frac{p+4}{3},
\end{equation}
and for every fixed $K\ge1$, a middle state with $\kappa\ge K$ satisfies
\begin{equation}\label{eq:middle-shell-K}
 R\le\frac{p+4}{4K-1}.
\end{equation}
\end{proposition}

\begin{proof}
Using $p=4hrs-R$ and $r+s=\kappa R$, one has the exact identity
\begin{equation}\label{eq:middle-bound-factorization}
 p+4-(4\kappa-1)R
 =4(hrs-r-s+1)
 =4\bigl((h-1)rs+(r-1)(s-1)\bigr).
\end{equation}
Both terms on the last side are nonnegative, proving
\eqref{eq:middle-shell-bound}.  The two displayed consequences follow from
$\kappa\ge1$ and $\kappa\ge K$, respectively.  Equality in
\eqref{eq:middle-shell-bound} holds exactly when $h=1$ and
$(r-1)(s-1)=0$.
\end{proof}

\section{Exact conclusion for Erd\H{o}s--Straus}

For every prime $p\equiv1\pmod4$, the complete result can be written as the
following equality of existence statements:
\begin{equation}\label{eq:master}
\boxed{
\begin{aligned}
\ES(p)
\quad\Longleftrightarrow\quad
\exists\;&\varepsilon\in\{0,1\},\;
0\le j<\frac{p-1}{4},\;
\ell_0,\ell_1,u,a,v\in\Z,\\
& T=4\ell_0-1>0,\quad
 a=\frac{p+4j+3}{4},\quad
 u=\ell_1+jT>0,\quad v>0,\\
&12\ell_0-4\ell_1-3=p^\varepsilon,\qquad
 uv=a^2.
\end{aligned}}
\end{equation}
The two equations on the last line are exactly
\[
 \det\mathsf R(\ell_0,\ell_1,1)=-p^\varepsilon,
 \qquad
 \det\Gamma(a,u)=0.
\]
The first is the terminal determinant of the $(3,4,\infty)$ sphere
construction.  The second is the positive rank-one Koide/$S^6$ coefficient
boundary.  The coupling equations on the preceding line preserve the prime,
first-denominator index, orientation, and divisor information.

The sphere construction therefore supplies an exact terminal-topological
interpretation of each selector state and the Koide coefficient form supplies
an exact rank-one completion of its divisor gate.  What remains in the
Erd\H{o}s--Straus conjecture is precisely the assertion that the integral
incidence set in \eqref{eq:master} is nonempty for every relevant prime.  The
$S^6$ terminal algebra alone does not force $uv=a^2$; the Koide determinant
boundary alone does not force
$12\ell_0-4\ell_1-3=p^\varepsilon$.  Their coupled integral incidence is the
complete selector.

Finally, the positive defect-one Koide/Vieta shell studied elsewhere in the
source has
\[
 \Delta=3Q_K=-1,
 \qquad Q_K=-\frac13,
\]
whereas every Erd\H{o}s--Straus selector and shell matrix constructed here has
$Q_K=0$.  An invertible isometry of the quadratic space preserves the value of
$Q_K$, so the defect-one shell is not the Erd\H{o}s--Straus null shell.  The
proved linear counting theorem for that defect-one hyperboloid is therefore a
separate arithmetic result, not a count of Erd\H{o}s--Straus witnesses.

\section*{Reproducibility}

The accompanying exact-integer program verifies every displayed map,
congruence, inverse, determinant, denominator identity, quarter-coordinate
formula, channel overlap, and double-certificate reconstruction for every
prime $p\equiv1\pmod4$ through $2000$.  It checks $147$ primes and $6548$
tagged selector states.  They give exactly $4963$ sorted denominator triples,
in agreement with \eqref{eq:solution-count}; all $6548$ double certificates
are distinct, while $391$ terminal-matrix collision classes demonstrate the
information lost by the terminal projection alone.  The bounded computation
is a verification receipt, not a proof of universal selector nonemptiness.

\begin{thebibliography}{9}

\bibitem{AlpogeFableS6}
L. Alp\"oge, with the disclosed Fable/Claude contribution,
\emph{The $(3,4,\infty)$ modular family of $2$-tori, completed at its three
special points, is a complex structure on $S^6$}, together with the current
canonical reconstruction and integrated audit, Zenodo record 22238779 (2026).

\bibitem{ClankersES}
The Clankers,
\emph{Erd\H{o}s--Straus Project Archive: Residual Divisor Shells, Quotient
Surpluses, and Finite Verification}, Zenodo record 22209908 (2026), and the
active research continuation supplied with this note.

\bibitem{Koide1983}
Y. Koide,
\emph{New View of Quark and Lepton Mass Hierarchy},
Physical Review D \textbf{28} (1983), 252--254;
erratum, Physical Review D \textbf{29} (1984), 1544.

\bibitem{Kocik2012}
J. Kocik,
\emph{The Koide Lepton Mass Formula and Geometry of Circles},
arXiv:1201.2067 (2012).

\end{thebibliography}

\end{document}
